\documentclass[11pt]{amsart}
\usepackage[T1]{fontenc}
\usepackage{amsmath,amssymb,amsthm}
\usepackage[margin=1.1in]{geometry}
\usepackage{booktabs}
\usepackage{enumitem}
\usepackage{hyperref}

\newtheorem{theorem}{Theorem}[section]
\newtheorem{lemma}[theorem]{Lemma}
\newtheorem{proposition}[theorem]{Proposition}
\newtheorem{corollary}[theorem]{Corollary}
\newtheorem{conjecture}[theorem]{Conjecture}
\theoremstyle{definition}
\newtheorem{definition}[theorem]{Definition}
\newtheorem{remark}[theorem]{Remark}
\newtheorem{example}[theorem]{Example}

\setlength{\emergencystretch}{3em}
\newcommand{\cD}{\mathcal D}
\newcommand{\lean}[1]{\quad{\footnotesize[Lean: \texttt{#1}]}}
\newcommand{\av}{\mathrm{av}}
\newcommand{\nev}{\mathrm{ne}}
\newcommand{\te}{\mathrm{te}}

\title[Constrained Lov\'asz decompositions]{Constrained Lov\'asz decompositions:\\
a constraint language, its compatibility boundary, and Fan's lemma with end paths}
\author{Justin Williams}
\date{November 2026 (working draft)}

\begin{document}
\begin{abstract}
Gallai's path-decomposition conjecture has been reduced, in this project, to several open
statements about \emph{Lov\'asz decompositions}: decompositions of a graph with at most one
vertex $g$ of positive even degree into paths, in which every vertex $v$ is an end of exactly
$\deg v\bmod 2$ paths. Each statement prescribes properties of the paths $P_v$ ending at a few
vertices. We set up a common constraint language for them (avoidance ``$P_c$ avoids $a$'',
distinctness ``$P_u\ne P_v$'', terminal edges), and determine, by an exact search over
\emph{all} Lov\'asz decompositions of every graph of order at most nine, which constraint systems
are always satisfiable. The answer is sharp in the data: a system is satisfiable whenever its
constraint digraph is \emph{bipartite} (no avoided vertex is itself constrained), has
\emph{in-degree at most one} (every vertex is avoided by at most one end path), and each avoid set
misses $N[c]\cup N(g)$ for its own tail $c$, provided $\deg g\ne2$; with $\deg g=2$ one needs in
addition that no avoided vertex is adjacent to any constrained vertex. We call this the
constrained Lov\'asz conjecture CLT-B. It contains the local form of $L_1^*$ (hence $L_\infty$ and
EE) and $L_N$. Each of its conditions is necessary: we give minimal counterexamples, with hand
proofs, to every relaxation. The dense $T(6)$ base of the t6-dense paper lies outside the
language's local region, yet in the dense residual \emph{every} labelling has a base.

On the proof side we formalise Fan's lemma with end paths (a Fan transfer keeps every end that is
not a chain vertex on a path with the same vertex set, plus the target), and derive that Lov\'asz's
Case~I carries an arbitrary constraint system across a free target, with no admissibility of the
child. We locate exactly where Lov\'asz's induction for CLT-B stops: at the multi-constraint
version of the blocked configuration of $L_1^*$, and at adjacencies between tails.
\end{abstract}
\maketitle

\section{Introduction}\label{sec:intro}

Let $G$ be a finite simple graph with at most one vertex of positive even degree; if it exists we
call it $g$. A \emph{Lov\'asz decomposition} of $G$ is a decomposition of $E(G)$ into paths in which
every vertex $v$ is an end of exactly $\deg v\bmod2$ paths~\cite{Lov68}. Every odd vertex $v$ is
then an end of exactly one path, which we call $P_v$ (the \emph{end path} of $v$).

The project has reduced Gallai's conjecture, and several of its known hard special cases, to
statements of the form \emph{some Lov\'asz decomposition has prescribed properties of a few end
paths}. The open ones at the time of writing are (synthesis \S23; connections \S\S V--AD):
\begin{itemize}
\item the \emph{local form of $L_1^*$}: $V(P_c)\subseteq N[c]\cup N(g)$ for a prescribed odd $c$
\cite{L1S}; with the tight transfer conjecture TR it is equivalent to the \emph{some-target}
conjecture ST of~\cite{L1I};
\item $L_N$ and MSEP: in an all-odd graph, for a matching of non-adjacent pairs $(c_i,a_i)$, the
path $P_{c_i}$ avoids $a_i$ (resp.\ no path contains both, at distance $\ge3$) \cite{RD,MS};
\item the \emph{t6 base conjecture} \cite[Conj.~5.1]{T6}: five or six end-path conditions at six
vertices of a nearly complete $K_6$;
\item and, since this work began, the \emph{triangle base} and the disjunctive avoidance conjecture
DL~\cite{DT}, and the forest base~\cite{DT}.
\end{itemize}
The coordinator's question was whether one \emph{constrained Lov\'asz theorem} lies behind all of
them: a constraint language, a locality condition and a compatibility condition under which a
decomposition satisfying all constraints always exists.

\subsection*{Results}
\begin{enumerate}[label=(\arabic*)]
\item \emph{Language} (Section~\ref{sec:lang}). Atoms $\av(c,a)$, $\nev(u,v)$, $\te(v,w)$; the
leaf reductions between them; every open conjecture above written as a constraint system
(Table~\ref{tab:targets}).
\item \emph{The boundary} (Sections~\ref{sec:conj}, \ref{sec:sharp}). An exact oracle that
enumerates \emph{all} Lov\'asz decompositions of every graph of order $\le9$ (with at most one even
vertex) and computes all minimal unsatisfiable systems of a given shape. The surviving statement is
Conjecture~\ref{conj:cltb} (CLT-B): bipartite, in-degree $\le1$, per-tail locality,
$\deg g\ne2$; and its cross-local form CLT-B1 for all $g$. No counterexample among all graphs
of order $\le 9$ that were enumerated (Table~\ref{tab:data}) and $5{,}400$ random maximal
instances of order $11$--$16$. Every relaxation fails, and its minimal counterexamples
(Section~\ref{sec:sharp}) are on $3$--$7$ vertices.
\item \emph{Fan's lemma with end paths} (Theorem~\ref{thm:fanend}, Lean). This is Lemma~4.1
of~\cite{T6}, formalised; and its consequence for constraint systems, Lov\'asz's Case~I across a
free target (Theorem~\ref{thm:caseI}, Lean).
\item \emph{The induction} (Section~\ref{sec:ind}). The exact list of Lov\'asz moves that carry a
CLT-B system, with their conditions; the residual configurations; for one constraint they are the
blocked instances of~\cite{L1S}.
\item \emph{The t6 base} (Section~\ref{sec:t6}). In $30$ residual graphs of all five dense shapes
every one of the $16{,}032$ admissible labellings has an (E7) append base. The exact search stalls
on these instances; a local search finds the bases at once.
\end{enumerate}

\section{The constraint language}\label{sec:lang}

Throughout, $G$ has at most one vertex $g$ of positive even degree, and ``decomposition'' means
Lov\'asz decomposition. For an odd vertex $v$, $P_v$ is its end path.

\begin{definition}\label{def:atoms}
An \emph{atom} is one of
\begin{itemize}
\item $\av(c,a)$, $c$ odd, $a\ne c$: $a\notin V(P_c)$ (\emph{$c$ avoids $a$}; $c$ is the
\emph{tail}, $a$ the \emph{head});
\item $\nev(u,v)$, $u\ne v$ odd: $P_u\ne P_v$;
\item $\te(v,w)$, $vw\in E$, $v$ odd: the edge of $P_v$ at $v$ is $vw$.
\end{itemize}
A \emph{system} is a set of atoms; it is \emph{satisfiable} in $G$ if some decomposition satisfies
all of them. Its \emph{constraint digraph} has an arc $c\to a$ for every $\av(c,a)$. We write
$\av(c,A)$ for $\{\av(c,a):a\in A\}$.
\end{definition}

\subsection*{Leaf reductions} Terminal edges and avoidance interconvert through leaf graphs
(route-descent \S V.4a, multi-sep \S W.3, prescribed-end \S Y.2):
\begin{itemize}
\item $\av(c,A)$ in $G$ $\iff$ in $G+c\ell$ (a new leaf at $c$, which makes $c$ even), no path
contains $\ell$ and a vertex of $A$. Only one path contains $\ell$, so this is a vertex-set
property, and Fan's lemma with vertex sets preserves it.
\item $\te(v,w)$ in $G$ $\iff$ in $G-vw+w\ell$ the path through $\ell$ avoids $v$; here $v$ becomes
even.
\end{itemize}
So every system is a set of separations of leaves in a graph with several even vertices, one per
leaf; and conversely. The number of even vertices is the obstruction: Lov\'asz's theorem and the
separation theorem need at most one. This is why the language is stated in $G$, with end paths.

\begin{table}[t]\footnotesize
\begin{tabular}{@{}lll@{}}
\toprule
statement & system & where it sits\\
\midrule
$L_1^*$ local form \cite{L1S} & $\av(c,A)$, $A=V\setminus(N[c]\cup N(g))$ & one tail: CLT-B1 with $k=1$\\
$L_\infty$ & $\av(c,V\setminus N[c])$, $G$ all odd & CLT-B1, $k=1$\\
EE \cite{PE} & $\te(u,g)$, $u\in N(g)$ & via $L_1^*$ (\texttt{ee\_of\_L1StarGen})\\
ST \cite{L1I} & $\av(c,A)$ in $G-gx$, some $x\in N(g)$ & child is non-local ($A\cap N(x)\ne\emptyset$)\\
$L_N$ \cite{RD} & $\av(c_i,a_i)$, a matching, $G$ all odd & CLT-B (bipartite, in-degree 1)\\
MSEP \cite{MS} & separations, disjoint pairs at distance $\ge3$ & vertex-set form; 2-cycles\\
t6 base (E7) \cite{T6} & $z\to x_1\to y_2\to x_2\to y_3\to x_3{=}g$, $\nev$ on $\{z,y_2,y_3\}$ & a chain; non-local\\
triangle base, DL \cite{DT} & for each pair, $\av(a_i,b_i)$ or $\av(b_i,a_i)$ & free orientation; non-local\\
forest base \cite{DT} & $\av(v,p(v))$ (child avoids parent), $\nev$ on siblings & in-forest; in-degree $>1$\\
RR \cite{RD} & $\av(c_p,V(R_p)-c_p)$, $\nev$ for crossing routes & Gallai-strength\\
\bottomrule
\end{tabular}
\caption{The open conjectures as constraint systems.}\label{tab:targets}
\end{table}

\section{The constrained Lov\'asz conjecture}\label{sec:conj}

\begin{definition}\label{def:blocal}
Let $c_1,\dots,c_k$ be distinct odd vertices (the \emph{tails}) and $A_1,\dots,A_k$ sets of
vertices. The system $\bigcup_i\av(c_i,A_i)$ is \emph{locally bipartite of in-degree one}
(\emph{B-local}) if
\begin{enumerate}[label=(B\arabic*)]
\item\label{B1} the $A_i$ are pairwise disjoint (every vertex is avoided by at most one end path);
\item\label{B2} no tail is avoided: $A_i\cap\{c_1,\dots,c_k\}=\emptyset$;
\item\label{B3} $A_i\cap(N[c_i]\cup N(g))=\emptyset$ for every $i$ (\emph{locality}; $g\in A_i$ is
allowed).
\end{enumerate}
It is \emph{cross-local} if moreover $A_i\cap N(c_j)=\emptyset$ for all $i,j$.
\end{definition}

\begin{conjecture}[CLT-B]\label{conj:cltb}
Let $G$ have at most one vertex $g$ of positive even degree, with $\deg g\ne2$. Then every B-local
system is satisfiable.\lean{CLTB}
\end{conjecture}

\begin{conjecture}[CLT-B1]\label{conj:cltb1}
Let $G$ have at most one vertex of positive even degree. Then every cross-local B-local system is
satisfiable.\lean{CLTB1}
\end{conjecture}

For $k=1$ both are the local form of $L_1^*$ (CLT-B1 at every $\deg g$;
\texttt{l1StarGen\_of\_cltB1}), hence $L_\infty$ and EE (\texttt{lInfty\_of\_cltB1},
\texttt{ee\_of\_cltB1}). For a matching of non-adjacent pairs in an all-odd graph CLT-B is $L_N$
(\texttt{lN\_of\_cltB}), which gives Gallai for every graph whose even vertices induce a subgraph
with a perfect matching~\cite{RD}. The degree-two exception is genuine
(Example~\ref{ex:deg2}); a vertex $g$ of degree two is a subdivided edge, and suppressing it
creates a parallel edge when its neighbours are adjacent, the multigraph regime where PTM also
fails~\cite{TM}.

\begin{table}[t]\footnotesize
\begin{tabular}{@{}lrrrrr@{}}
\toprule
$n$ & graphs & CLT-B1 & CLT-B ($\deg g\ne2$) & B-local, $\deg g=2$ & + 2-cycles at dist.\ $\ge3$\\
\midrule
4 & 2 & 0 & 0 & -- & 0\\
5 & 7 & 0 & 0 & 0 & 0\\
6 & 13 & 0 & 0 & -- & 0\\
7 & 130 & 0 & 0 & 2 & 0\\
8 & 224 & 0 & 0 & -- & 0\\
9 & $10{,}719^\dagger$ & 0 & 0 & 42 & 0\\
\bottomrule
\end{tabular}
\caption{Number of connected graphs with at most one even vertex having an unsatisfiable system
of the given class, exact over all systems with up to $8$--$10$ atoms (larger systems in a graph of
this order do not exist: a B-local system has at most $n-1$ atoms). $^\dagger$The $10{,}719$ of the $11{,}936$ graphs of order
$9$ enumerated so far; the others are the densest, see E30.}\label{tab:data}
\end{table}

\subsection*{Evidence}
\begin{itemize}
\item Exact, Table~\ref{tab:data}. For every graph, \texttt{enum.c} lists all Lov\'asz
decompositions and keeps the inclusion-minimal \emph{violation sets} (atoms violated by one
decomposition); a system is unsatisfiable iff it meets every one of them, and \texttt{trans.c}
searches for such hitting sets of the prescribed shape, smallest first.
\item Random, order $11$--$16$: $5{,}400$ maximal systems (every admissible vertex assigned to a
random tail; $1{,}800$ cross-local, $1{,}800$ cross-local with all tails' end paths distinct,
$1{,}800$ CLT-B with heads next to other tails), decided by the exact DFS of~\cite{T6}, the
undecided ones by local search (Section~\ref{sec:t6}); $0$ failures, $0$ left undecided for CLT-B,
$9$ left undecided (budget) for the cross-local forms.
\item The failures of B-local systems at $\deg g=2$ (2 at $n=7$, 42 at $n=9$) all have an avoided
vertex adjacent to another tail; under cross-locality there are none.
\end{itemize}

\section{The boundary: every relaxation fails}\label{sec:sharp}

Each condition of Definition~\ref{def:blocal} is necessary. The examples are minimal in the data
(no smaller graph or system fails). We write the graphs by adjacency lists.

\begin{example}[locality at $c$]\label{ex:Kn}
In $K_{2m}$ every Lov\'asz decomposition has $m$ paths and $m(2m-1)$ edges, so every path is
Hamiltonian, and $\av(c,a)$ fails for all $c\ne a$. So heads must avoid $N[c]$.
\end{example}

\begin{example}[locality at $g$]\label{ex:path}
In the path $c\,g\,a$ the only decomposition is the path itself, so $\av(c,a)$ fails; here
$a\in N(g)$. ($L_1$ with an even vertex fails in exactly this way, route-descent \S V.4a.)
\end{example}

\begin{example}[in-degree two]\label{ex:indeg}
Let $G$ be $K_4$ on $\{0,3,4,5\}$ minus $34$, with leaves $1$ at $3$ and $2$ at $4$ (all degrees odd,
$n=6$). The system $\av(1,\{0,5\})\cup\av(2,0)$ is cross-local and bipartite, and $0$ is
avoided twice. It is unsatisfiable. Indeed $P_1$ must be $13$ (the other neighbours of $3$ are
$0,5$), so $3$ has no further end and $0\,3\,5$ lies inside one path. $P_2$ avoids $0$ and cannot
end at $3$, so $P_2\in\{24,\,245\}$. If $P_2=24$, the remaining five edges must form one path
with ends $0,5$ through $3$ and $4$, impossible on four vertices. If $P_2=245$, the remaining
edges $03,04,05,35$ must form one path from $4$ to $0$, again impossible. Every proper subsystem is
satisfiable.
\end{example}

\begin{example}[chains: a head that is a tail]\label{ex:net}
Let $G$ be the net: a triangle $0\,4\,5$ with leaves $3$ at $0$, $1$ at $4$, $2$ at $5$ (all
odd, three paths). The system $\av(2,0),\av(0,1),\av(1,5),\av(3,4)$ is local, of in-degree one and
acyclic; its only defect is the chain $2\to0\to1\to5$. It is unsatisfiable. The end paths of the
three leaves are distinct: $P_1=P_3$ would be $1\,4\,0\,3$, containing $4$; $P_3=P_2$ would be
$3\,0\,5\,2$, containing $0$; $P_2=P_1$ would contain $5$. So the leaves are paired with
$\{0,4,5\}$, as $(1,4),(3,0),(2,5)$ or as $(1,0),(3,5),(2,4)$. The first leaves the triangle
uncovered. The second forces $P_1=1\,4\,0$, $P_3=3\,0\,5$, $P_2=2\,5\,4$, and then $P_0=P_1\ni1$.
\end{example}

\begin{example}[cycles]\label{ex:cyc}
In $K_{1,3}$ with leaves $a,b,c$ the directed triangle $\av(a,b),\av(b,c),\av(c,a)$ fails: the
two paths pair the four ends $a,b,c,z$, and each of the three pairings puts some leaf on a path with
the vertex it must avoid. The same holds for $\av(a,\{b,c\}),\av(b,c)$, a transitive triangle.
Simultaneous confinement $V(P_a)\subseteq N[a]$, $V(P_b)\subseteq N[b]$ is the system
$\av(a,\{b,c\}),\av(b,\{a,c\})$ and fails for the same reason (l1star \S AA.3). Two-cycles at
distance two fail as well: in the all-odd graph \texttt{EErw} with $N(0)=N(1)=\{3,4,5\}$,
$N(2)=\{5\}$, $N(3)=N(4)=\{0,1,5\}$, $N(5)=\{0,1,2,3,4\}$, the system
$\av(0,1),\av(1,0),\av(3,4),\av(4,3)$ (two 2-cycles between twins at distance two) fails
(exhaustive); this is the known failure of separation at distance two.
\end{example}

\begin{example}[cross-locality at $\deg g=2$]\label{ex:deg2}
Graph \texttt{FCRV?} ($n=7$): $N(0)=\{3,5,6\}$, $N(1)=\{4,5,6\}$, $N(2)=\{6\}$, $N(3)=\{0,5\}$,
$N(4)=\{1\}$, $N(5)=\{0,1,3\}$, $N(6)=\{0,1,2\}$; $g=3$ has degree two and lies on the triangle
$0\,3\,5$. The system $\av(0,1),\av(4,6)$ is B-local, but the head $6$ is adjacent to the tail $0$.
It is unsatisfiable (exhaustive). All $44$ failures of B-local systems at orders $7$ and $9$ have
$\deg g=2$ and such a cross adjacency.
\end{example}

\begin{remark}[distinctness and non-local atoms]
Distinctness atoms need their own compatibility: $\nev(u,v)$ fails whenever $u,v$ are the only odd
vertices, and adding $\nev$ between arbitrary non-avoided vertices makes every graph of order
$\le 8$ fail. Non-local avoidance atoms fail \emph{jointly} even when each is satisfiable alone: at
$n=7$ (exhaustive over all avoidance atoms; $n\le6$ has no such case), $18$ graphs have a matching-shaped system of individually satisfiable atoms that is
jointly unsatisfiable, e.g.\ \texttt{FCdcw} with $\av(2,0),\av(5,3),\av(1,4)$, all three heads in
$N(g)$, $\deg g=4$. This is the regime of ST and of DL; there the freedom (of the target $x$, of
the orientation) is used, and in the data of~\cite{DT} the free-orientation form DL has no
compatibility failure (every set of pairs whose disjunctions hold one at a time holds jointly).
\end{remark}

\section{Fan's lemma with end paths}\label{sec:fan}

Fan's lemma~\cite{Fan05} adds an edge $uv$ to a decomposition $\cD'$ of $G'=G-uv$ when
$\mathrm{ends}(v)$ exceeds the number of endless neighbours of $u$: an end moves from $v$ to $u$. Its
proof is a chain $v=x_0,x_1,\dots,x_k$ of re-routings followed by one append. A re-routing reverses
an initial segment of the path ending at $x_i$, trading the edge $x_{i+1}u$ for $x_iu$; it keeps the
vertex set and moves the end from $x_i$ to $x_{i+1}$. The append adds $u$ at the end $x_k$. Every
$x_i$ with $i\ge1$ precedes $u$ on a path, so it is a neighbour of $u$.

\begin{theorem}[Fan with end paths]\label{thm:fanend}
In the situation of Fan's lemma, every path $Q$ of the new decomposition $\cD$ of $G$ comes from one
path $Q'$ of $\cD'$: $V(Q)\subseteq V(Q')\cup\{u\}$, and every end $w$ of $Q$ with
$w\notin\{u,v\}\cup N_{G'}(u)$ is an end of $Q'$.\lean{fan\_single\_track}
\end{theorem}

\begin{proof}
For one re-routing, the re-routed path has the same vertex set and the same ends except that $x_i$
is replaced by $x_{i+1}$; all other paths are unchanged
(\texttt{reroute\_core\_track}). For the append, the extended path has the old vertex set plus $u$,
and its ends are $u$ and the old far end (\texttt{append\_move\_track}). These one-step statements
(\texttt{Tracks}) compose (\texttt{Tracks.trans}) along the chain (\texttt{fan\_exec\_track}),
excluding the chain vertices and $u$; the chain vertices after $v$ are neighbours of $u$
(\texttt{FanPre.adj\_of\_mem}).
\end{proof}

This is \cite[Lemma 4.1]{T6} (``an end at $v\notin N(u)\cup\{u\}$ stays on a path with the same
vertex set plus $u$''), which was paper only, in the stronger \emph{one-path} form: both ends of a
new path are traced to the same old path. That is what distinctness atoms need.

\begin{theorem}[Fan transfers carry systems]\label{thm:transfer}
If $\cD'$ satisfies a system $S$ (avoidance and distinctness atoms), no constrained vertex lies in
$\{u,v\}\cup N_{G'}(u)$, and $u$ is not avoided, then the decomposition of Fan's lemma satisfies
$S$.\lean{csys\_fan\_transfer}
\end{theorem}

\begin{proof}
Let $Q$ end at a tail $c$. By Theorem~\ref{thm:fanend}, $Q'$ ends at $c$, so $V(Q')\cap A=\emptyset$,
and $V(Q)\subseteq V(Q')\cup\{u\}$ with $u\notin A$. If $Q$ ends at $p$ and $q$ with $\nev(p,q)\in S$,
then $Q'$ ends at both, which $\cD'$ forbids.
\end{proof}

\begin{theorem}[Lov\'asz's Case I for systems]\label{thm:caseI}
Let $g$ be the even vertex, $u\in N(g)$, and let $S$ be a system such that no constrained vertex
lies in $N[u]\cup\{g\}$ and $u$ is not avoided. If $G-gu$ has a Lov\'asz decomposition satisfying
$S$, so does $G$.\lean{csys\_caseI}
\end{theorem}

\begin{proof}
In $G-gu$ only $u$ can be even, $g$ is odd with one end, and every neighbour of $u$ is odd with an
end, so Fan's lemma applies with $v=g$; the end counts become those of $G$. Apply
Theorem~\ref{thm:transfer}.
\end{proof}

No admissibility of the child is assumed. For one avoidance constraint this is
\texttt{avoid\_of\_target\_child} of~\cite{L1I}, which used a leaf graph; with several constraints
the leaf graph would have several even vertices, and the end-path form is what makes the
multi-constraint statement possible.

\subsection*{Packed and sparse systems} Two unconditional cases follow from the separation
theorems: if all tails and heads lie in a packed set (pairwise distance $\ge3$), every system of
avoidance and distinctness atoms among them is satisfiable (\texttt{csys\_good\_of\_packed}, from
\texttt{separation\_set}); if every vertex has at most one neighbour in the set of tails and heads,
every avoidance system with heads outside $N[c]$ is satisfiable (\texttt{csys\_good\_of\_sparse},
from \texttt{separation\_sparse\_set}), with no disjointness or cross-locality needed.

\section{Lov\'asz's induction for CLT-B}\label{sec:ind}

We run Lov\'asz's induction (measure $\sum\deg+[\text{no even vertex}]$) on a B-local instance
$(G,\{(c_i,A_i)\})$ and list every move that carries the system, with the exact condition under
which it applies and under which its child is again an instance. Transfers never need the child to
be an instance; only the call to the induction hypothesis does (as in~\cite{L1I}).

\begin{enumerate}[label=(M\arabic*)]
\item \emph{Free target} (Theorem~\ref{thm:caseI}). $u\in N(g)$, no tail in $N[u]$, $u\notin A$.
Child $(G-gu)$ with the same system; it is an instance iff $N(u)\cap A\subseteq\{g\}$ (locality at
the new even vertex $u$), and for CLT-B also $\deg u\ne3$ or cross-locality.
\item \emph{Tail at $g$} (Lean: \texttt{csys\_case\_adj}). $c_1\in N(g)$. Child $G-gc_1$, where $c_1$
is even; the tail $g$ replaces $c_1$ with avoid set $A_1\cup\{c_1\}$; append $gc_1$ to $P_g$. The transfer is unconditional (if
$P_g$ also ends at $c_j$, the new $P_{c_j}$ gains $c_1\notin A_j$). The child is an instance iff
$g\notin A$, $A_j\cap N(c_1)=\emptyset$ for $j\ne1$, and (cross-local form) $c_1$ is adjacent to
no other tail.
\item \emph{Tight target.} $u\in N(g)\cap N(c_1)$, $N(u)\cap A\subseteq\{g\}$, no other tail in
$N[u]$. Child $G-gu-c_1u$ with the tail $u$ replacing $c_1$ and avoid set $A_1\cup\{c_1\}$, as in
\texttt{avoid\_of\_target\_adj\_child}; the other tails are carried by Theorem~\ref{thm:transfer}.
Same instance conditions as (M2) with $u$ for $g$.
\item \emph{All odd, a tail of degree $\ge3$} (Lean: \texttt{csys\_case\_leaf}). $y\in N(c_1)$; the leaf graph $G-c_1y+y\ell$, where
$c_1$ is even and the leaf $\ell$ replaces $c_1$ with avoid set $A_1\cup\{c_1\}$; merge. The
transfer is unconditional; the child is an instance iff $A_j\cap N(c_1)=\emptyset$ for $j\ne1$
(and $c_1$ adjacent to no other tail, cross-local form).
\item \emph{All odd, a free vertex.} Lov\'asz's Case II at $x\notin N[\text{tails}]\cup A\cup N(A)$
of degree $\ge3$, by Theorem~\ref{thm:transfer}.
\end{enumerate}

\begin{proposition}\label{prop:residual}
Let $(G,\{(c_i,A_i)\})$ be a counterexample to CLT-B1 with $\sum\deg$ smallest, and
$A=\bigcup A_i$. Then none of (M1)--(M5) applies. In particular, if $g$ exists, then
\begin{enumerate}[label=(\alph*)]
\item every non-tail $u\in N(g)$ has a neighbour in $A\setminus\{g\}$, or two tails among its
neighbours, or exactly one tail $c$ among its neighbours and $c$ adjacent to another tail;
\item if some tail $c$ lies in $N(g)$, then $g\in A$ or $c$ is adjacent to another tail.
\end{enumerate}
\end{proposition}

\begin{proof}[Proof sketch]
Each move produces a smaller graph, and its transfer turns a good child into a good instance.
(M1) is Theorem~\ref{thm:caseI}. (M2) is the append of \texttt{l1\_case\_adj}~\cite{L1S}: the
other paths are unchanged, and a path that also ends at another tail $c_j$ gains only $c_1\notin
A_j$. (M3) is \texttt{avoid\_of\_target\_adj\_child}~\cite{L1I} (leaf graph at $u$ standing for
$c_1$, Fan's lemma towards $u$, merge); in the leaf graph the other tails are not in $N[u]\cup\{g\}$,
so Theorem~\ref{thm:transfer} carries their constraints, and the merge changes only the leaf path,
which avoids $c_1$. (M4) is the leaf merge of \texttt{l1\_case\_big}, which changes only the leaf
path. (M5) is Lov\'asz's Case~II with Theorem~\ref{thm:transfer}. The instance conditions stated in
each move are checked directly against Definition~\ref{def:blocal} and cross-locality; for
example in (M3) the new even vertex is $c_1$, and $A_j\cap N(c_1)=\emptyset$ by cross-locality,
while the new head $c_1$ is adjacent to no tail other than through deleted edges unless $c_1$ is
adjacent to another tail. If (a) failed for some $u$, (M1) or (M3) would apply at $u$; if (b)
failed, (M2) would apply. (M1), (M2) and (M4) are formalised (\texttt{csys\_caseI}, \texttt{csys\_case\_adj},
\texttt{csys\_case\_leaf}, for arbitrary further constraints); (M3) and (M5) for several tails are
not.
\end{proof}

For $k=1$ the residual is exactly the blocked configuration of~\cite{L1S}
(\texttt{l1starGen\_of\_blocked}), and the analysis of~\cite{L1I} applies: the blocked case splits
into ST and TR. For $k\ge2$ two new obstructions appear: tails adjacent to each other (which (M2),
(M4) turn into heads adjacent to tails, leaving the cross-local class), and heads of one tail next
to another tail (which leave the class at the new even vertex). Neither occurs when the tails are
pairwise at distance $\ge3$: then no vertex is adjacent to two tails, and the residual is the
blocked configuration of~\cite{L1S} for the union $A$ of the avoid sets (every non-tail neighbour of
$g$ has a neighbour in $A\setminus\{g\}$, and $g\in A$ if a tail is adjacent to $g$).

\section{The t6 base and the other non-local systems}\label{sec:t6}

The t6 base system (Table~\ref{tab:targets}) violates every condition of CLT-B: it is a chain, its
heads are adjacent to tails in the dense regime, and two heads lie in $N(x_3)$. The analysis
of~\cite[AD.4]{T6} shows the non-local conditions are binding for some labellings. Our data shows
that they are not binding for the \emph{existence} of a base in residual graphs:

\begin{itemize}
\item On $30$ residual graphs (\texttt{gen\_resid.py}, $|O|=8$, six per shape $C_4,K_3,P_3,K_2,\emptyset$),
\emph{every} labelling with $x_1y_2,x_2y_3\in E(B)$ ($368$--$720$ per graph, $16{,}032$ in all)
has an (E7) append base; every witness checked. The same holds on $15$ residual graphs with
$|O|=10$ ($8{,}016$ labellings) and on $6$ sparse $C_4$ graphs of \texttt{gen\_core.py}
($2{,}208$ labellings): $26{,}256$ labellings, no exception. (In small graphs there are exceptions:
at order $9$, $16$ of $14{,}144$ labellings have no (E7) base~\cite[E28]{T6}.)
\item The exact DFS (\texttt{lsearch.c}, budget $3\cdot10^5$ nodes, five vertex orders) decides
$9{,}188$ of them (all positively); on whole graphs it stalls even at $10^7$ nodes, although an
unconstrained decomposition of the same graph is found instantly. A local search
(\texttt{lsls.py}: delete three paths, one of them violating, and re-decompose their union exactly
with the constraints on its ends) finds every remaining base in under a second.
\end{itemize}
So in the dense regime the t6 base is a universal statement in the data, and the difficulty is
the proof, not the choice of labelling. The same local search resolves the budget-undecided roots
of the random CLT-B tests.

For ST the child $(G-gx,c,A)$ is non-local at the new even vertex $x$. The universal form
``$L_1^*$ at $c$ together with EE at every $x\in N(g)$'' fails ($39$ of $2{,}210$ instances at
$n=7$), always with $N(x)\cap A\ne\emptyset$; the existential form (some $x$), which implies ST,
never fails to $n=7$. This is the L1E statement of~\cite{L1I}.

\section{Computations}\label{sec:comp}
\emph{This section is computational; no proof above depends on it.} Scripts are in
\texttt{scripts/}; the log is E30 in \texttt{research/experiments.md}.
\begin{itemize}
\item \texttt{enum.c}: all Lov\'asz decompositions by the DFS of~\cite{STL,T6} (paths grown from a
deterministic vertex with an end, so each decomposition once; dead states memoised), with the
inclusion-minimal violation sets over the atom universe (local or all avoidance atoms, distinctness
atoms; separation atoms for graphs with packed even sets).
\item \texttt{trans.c}: smallest hitting sets of the violation family whose atoms form a system of a
given shape (iterative deepening, branching on the unhit set with fewest admissible atoms). Shapes:
star, matching, bipartite, in-degree one, functional, acyclic, 2-cycles at distance $\ge3$, with
degree filters.
\item \texttt{rand\_clt.py}: random connected graphs with parity repaired, maximal systems, exact
DFS then local search.
\item \texttt{t6all.py}: all labellings of the t6 base on residual graphs.
\end{itemize}

\section{The formalisation}\label{sec:lean}
\texttt{lean/GallaiLean/ConstrainedLovasz.lean}, sorry-free, standard axioms, full build green:
\texttt{Tracks}, \texttt{append\_move\_track}, \texttt{reroute\_core\_track},
\texttt{reroute\_move\_track}, \texttt{fan\_exec\_track}, \texttt{fan\_single\_track}
(Theorem~\ref{thm:fanend}); \texttt{CSys}, \texttt{CSys.Sat}, \texttt{CSys.Good},
\texttt{CSys.sat\_of\_tracks}, \texttt{csys\_fan\_transfer} (Theorem~\ref{thm:transfer}),
\texttt{csys\_caseI} (Theorem~\ref{thm:caseI}); \texttt{BLocal}, \texttt{CrossLocal},
\texttt{CLTB}, \texttt{CLTB1}, \texttt{LN} and \texttt{l1StarGen\_of\_cltB1},
\texttt{lInfty\_of\_cltB1}, \texttt{ee\_of\_cltB1}, \texttt{lN\_of\_cltB};
\texttt{csys\_good\_of\_packed}, \texttt{csys\_good\_of\_sparse}. \texttt{csys\_case\_adj}, \texttt{CSys.liftLeaf}, \texttt{csys\_case\_leaf} (moves (M2), (M4) for
several tails). Not formalised: the moves (M3), (M5) for several tails, Proposition~\ref{prop:residual}, and the examples of
Section~\ref{sec:sharp} (checked by hand or exhaustively).

\begin{thebibliography}{99}
\bibitem{Fan05} G.~Fan, Path decompositions and Gallai's conjecture, J. Combin. Theory Ser. B 93
(2005) 117--125.
\bibitem{Lov68} L.~Lov\'asz, On covering of graphs, in: Theory of Graphs (Proc. Colloq. Tihany,
1966), Academic Press, 1968, 231--236.
\bibitem{RD} This project, \texttt{manuscripts/route-descent} (route descent, $L_1$, $L_N$, RR).
\bibitem{MS} This project, \texttt{manuscripts/multi-sep} (multi-pair separation).
\bibitem{PE} This project, \texttt{manuscripts/prescribed-end} (EE).
\bibitem{STL} This project, \texttt{manuscripts/stl}.
\bibitem{TM} This project, \texttt{manuscripts/terminal-matching} (PTM).
\bibitem{L1S} This project, \texttt{manuscripts/l1star} ($L_1^*$).
\bibitem{L1I} This project, \texttt{manuscripts/l1-induction} (TR, ST).
\bibitem{T6} This project, \texttt{manuscripts/t6-dense} (the t6 base conjecture).
\bibitem{DT} This project, \texttt{manuscripts/direct-test1} (the direct lift, triangle and forest
bases, DL).
\end{thebibliography}
\end{document}
